% Text-native transcription of geography-sl-hl-2017-en.pdf.
% Optimized for AI consumption: source wording is preserved, decorative layout
% is omitted, and text embedded in the front-page programme graphic is captured
% explicitly in prose below.
\documentclass[11pt,a4paper]{article}

\usepackage[utf8]{inputenc}
\usepackage[T1]{fontenc}
\usepackage{lmodern}
\usepackage[margin=1in]{geometry}
\usepackage{array}
\usepackage{booktabs}
\usepackage{enumitem}
\usepackage{hyperref}
\usepackage{tabularx}

\hypersetup{
    colorlinks=true,
    linkcolor=black,
    urlcolor=blue,
    pdftitle={IB DP Subject Brief -- Geography SL \& HL},
    pdfauthor={International Baccalaureate Organization}
}

\setlength{\parindent}{0pt}
\setlength{\parskip}{0.6em}
\setlist[itemize]{leftmargin=1.5em,itemsep=2pt,topsep=2pt}
\setlist[enumerate]{leftmargin=1.8em,itemsep=3pt,topsep=3pt}

\newcolumntype{Y}{>{\raggedright\arraybackslash}X}

\begin{document}

\begin{center}
{\Large\textbf{International Baccalaureate}}\\[2pt]
{\Large\textbf{Diploma Programme Subject Brief}}\\[6pt]
{\large Individuals and societies: Geography}\\[4pt]
First assessments 2019\\[6pt]
\copyright{} International Baccalaureate Organization 2017
\end{center}

The IB Diploma Programme (DP) is a rigorous, academically challenging and balanced
programme of education designed to prepare students aged 16 to 19 for success at
university and life beyond. The DP aims to encourage students to be knowledgeable,
inquiring, caring and compassionate, and to develop intercultural understanding,
open-mindedness and the attitudes necessary to respect and evaluate a range of
viewpoints. Approaches to teaching and learning (ATL) are deliberate strategies,
skills and attitudes that permeate the teaching and learning environment. In the
DP students develop skills from five ATL categories: thinking, research, social,
self-management and communication.

To ensure both breadth and depth of knowledge and understanding, students must
choose at least one subject from five groups: 1) their best language, 2) additional
language(s), 3) social sciences, 4) sciences, and 5) mathematics. Students may choose
either an arts subject from group 6, or a second subject from groups 1 to 5. At least
three and not more than four subjects are taken at higher level (240 recommended
teaching hours), while the remaining are taken at standard level (150 recommended
teaching hours). In addition, three core elements---the extended essay, theory of
knowledge and creativity, activity, service---are compulsory and central to the
philosophy of the programme.

These IB DP subject briefs illustrate the following key course components.
\begin{enumerate}[label=\Roman*.]
    \item Course description and aims
    \item Curriculum model overview
    \item Assessment model
    \item Sample questions
\end{enumerate}

\subsection*{Front-page programme graphic labels in the source PDF}
The source PDF includes a central IB Diploma Programme graphic. Its text labels are:
\begin{itemize}
    \item Outer programme labels: IB Diploma Programme; International-mindedness.
    \item Subject-group labels: Studies in Language and Literature; Language Acquisition;
    Individuals and Societies; Sciences; Mathematics; The Arts.
    \item Core labels: Theory of Knowledge; Extended Essay; Creativity, Activity, Service.
    \item Inner labels: Approaches to Teaching; Approaches to Learning.
\end{itemize}

\section*{I. Course description and aims}

Geography is a dynamic subject firmly grounded in the real world, and focuses on
the interactions between individuals, societies and physical processes in both
time and space. It seeks to identify trends and patterns in these interactions.
It also investigates the way in which people adapt and respond to change, and
evaluates actual and possible management strategies associated with such change.
Geography describes and helps to explain the similarities and differences between
different places, on a variety of scales and from different perspectives.

Geography as a subject is distinctive in its spatial dimension and occupies a
middle ground between social or human sciences and natural sciences. The course
integrates physical, environmental and human geography, and students acquire
elements of both socio-economic and scientific methodologies. Geography takes
advantage of its position to examine relevant concepts and ideas from a wide
variety of disciplines, helping students develop life skills and have an
appreciation of, and a respect for, alternative approaches, viewpoints and ideas.

Students at both SL and HL are presented with a common core and optional
geographic themes. HL students also study the HL core extension. Although the
skills and activity of studying geography are common to all students, HL students
are required to acquire a further body of knowledge, to demonstrate critical
evaluation and to further synthesize the concepts in the HL extension.

The aims of the geography course at SL and HL are to enable students to:
\begin{itemize}
    \item develop an understanding of the dynamic interrelationships between
    people, places, spaces and the environment at different scales
    \item develop a critical awareness and consider complexity thinking in the
    context of the nexus of geographic issues, including:
    \begin{itemize}[label=$\circ$,leftmargin=1.5em]
        \item acquiring an in-depth understanding of how geographic issues,
        or wicked problems, have been shaped by powerful human and
        physical processes
        \item synthesizing diverse geographic knowledge in order to form
        viewpoints about how these issues could be resolved.
    \end{itemize}
    \item understand and evaluate the need for planning and sustainable
    development through the management of resources at varying scales.
\end{itemize}

\section*{II. Curriculum model overview}

\begin{tabularx}{\textwidth}{>{\raggedright\arraybackslash}p{0.74\textwidth} >{\centering\arraybackslash}p{0.09\textwidth} >{\centering\arraybackslash}p{0.09\textwidth}}
\toprule
\textbf{Syllabus component} & \multicolumn{2}{c}{\textbf{Teaching hours}} \\
\cmidrule(lr){2-3}
 & \textbf{SL} & \textbf{HL} \\
\midrule
\textbf{Geographic themes---seven options}\\
SL---two options; HL---three options
\begin{itemize}[leftmargin=1.2em,itemsep=1pt,topsep=1pt]
    \item Freshwater
    \item Oceans and coastal margins
    \item Extreme environments
    \item Geophysical hazards
    \item Leisure, tourism and sport
    \item Food and health
    \item Urban environments
\end{itemize}
& 60 & 90 \\
\addlinespace[0.4em]
\textbf{SL and HL core}\\
Geographic perspectives---global change
\begin{itemize}[leftmargin=1.2em,itemsep=1pt,topsep=1pt]
    \item Population distribution---changing population
    \item Global climate---vulnerability and resilience
    \item Global resource consumption and security
\end{itemize}
& 70 & 70 \\
\addlinespace[0.4em]
\textbf{HL only}\\
Geographic perspectives---global interactions
\begin{itemize}[leftmargin=1.2em,itemsep=1pt,topsep=1pt]
    \item Power, places and networks
    \item Human development and diversity
    \item Global risks and resilience
\end{itemize}
&  & 60 \\
\addlinespace[0.4em]
\textbf{Internal assessment}\\
SL and HL Fieldwork\\
Fieldwork, leading to one written report based on a fieldwork question,
information collection and analysis with evaluation
& 20 & 20 \\
\midrule
\textbf{Total teaching hours} & \textbf{150} & \textbf{240} \\
\bottomrule
\end{tabularx}

\section*{III. Assessment model}

There are four assessment objectives (AOs) for the SL and HL geography course.
Having followed the course at SL or HL, students will be expected to do the
following:

\begin{enumerate}
    \item \textbf{Demonstrate knowledge and understanding of specified content}
    \begin{itemize}
        \item the core theme---global change
        \item two optional themes at SL and three optional themes at HL
        \item at HL, the HL extension---global interactions
        \item in internal assessment, a specific geographic research topic.
    \end{itemize}
    \item \textbf{Demonstrate application and analysis of knowledge and understanding}
    \begin{itemize}
        \item apply and analyse geographic concepts and theories
        \item identify and interpret geographic patterns and processes in
        unfamiliar information, data and cartographic material
        \item demonstrate the extent to which theories and concepts are recognized
        and understood in particular contexts.
    \end{itemize}
    \item \textbf{Demonstrate synthesis and evaluation}
    \begin{itemize}
        \item examine and evaluate geographic concepts, theories and perceptions
        \item use geographic concepts and examples to formulate and present
        an argument
        \item evaluate materials using methodology appropriate for geographic
        fieldwork
        \item at HL only, demonstrate synthesis and evaluation of the HL
        extension---global interactions.
    \end{itemize}
    \item \textbf{Select, use and apply a variety of appropriate skills and techniques}
    \begin{itemize}
        \item select, use and apply:
        \begin{itemize}[label=$\circ$,leftmargin=1.5em]
            \item prescribed geographic skills in appropriate contexts
            \item techniques and skills appropriate to a geographic research
            question.
        \end{itemize}
        \item produce well-structured written material, using appropriate
        terminology.
    \end{itemize}
\end{enumerate}

\subsection*{Assessment at a glance}

\begin{tabularx}{\textwidth}{>{\raggedright\arraybackslash}p{0.13\textwidth} >{\raggedright\arraybackslash}Y >{\centering\arraybackslash}p{0.08\textwidth} >{\centering\arraybackslash}p{0.08\textwidth} >{\centering\arraybackslash}p{0.08\textwidth} >{\centering\arraybackslash}p{0.08\textwidth}}
\toprule
\textbf{Type of assessment} & \textbf{Format of assessment} & \multicolumn{2}{c}{\textbf{Time (hours)}} & \multicolumn{2}{c}{\textbf{Weighting of final grade (\%)}} \\
\cmidrule(lr){3-4}\cmidrule(lr){5-6}
 & & \textbf{SL} & \textbf{HL} & \textbf{SL} & \textbf{HL} \\
\midrule
\textbf{External} &  & 2.75 & 4.5 & 75 & 80 \\
Paper 1 & Each option has a structured question and one extended answer question from a choice of two. & 1.5 & 2.25 & 35 & 35 \\
Paper 2 & Three structured questions, based on each SL/HL core unit. Infographic or visual stimulus, with structured questions. One extended answer question from a choice of two. & 1.25 & 1.25 & 40 & 25 \\
Paper 3 & Choice of three extended answer questions, with two parts, based on each HL core extension unit. &  & 1 &  & 20 \\
\addlinespace[0.2em]
\textbf{Internal} &  & 20 & 20 & 25 & 20 \\
Fieldwork & One written report based on a fieldwork question from any suitable syllabus topic, information collection and analysis with evaluation. & 20 & 20 & 25 & 20 \\
\bottomrule
\end{tabularx}

\section*{IV. Sample questions}

\begin{itemize}
    \item Examine the role of plate margin type in determining the severity of
    volcanic hazards.
    \item Evaluate the success of attempts to predict tectonic hazard event
    and their possible impacts.
    \item Evaluate the role of agribusiness and new technologies in increasing
    world food supply.
    \item Examine the relationship between food security and health.
    \item Using examples, analyse how technological developments can threaten
    the security of states.
    \item To what extent does a global culture exist?
\end{itemize}

\subsection*{About the IB}

For nearly 50 years, the IB has built a reputation for high-quality, challenging
programmes of education that develop internationally minded young people who are
well prepared for the challenges of life in the 21st century and are able to
contribute to creating a better, more peaceful world.

For further information on the IB Diploma Programme, and a complete list of DP
subject briefs, visit: \url{http://www.ibo.org/diploma/}.

Complete subject guides can be accessed through the IB online curriculum centre
(OCC) or purchased through the IB store: \url{http://store.ibo.org}.

For more on how the DP prepares students for success at university, visit:
\url{www.ibo.org/recognition} or email:
\href{mailto:recognition@ibo.org}{recognition@ibo.org}.

\bigskip

\small
International Baccalaureate\textregistered{} \,|\, Baccalaur\'eat International\textregistered{} \,|\, Bachillerato Internacional\textregistered{}

\end{document}
